\chapter{Содержание работы}

\textbf{Цель работы:} изучение венгерского метода решения задачи о назначениях.

\begin{enumerate}
    \item реализовать венгерский метод решения задачи о назначениях в виде программы на ЭВМ;
    \item провести решение задачи с матрицей стоимостей, заданной в индивидуальном варианте, рассмотрев два случая:
    \begin{itemize}
        \item [a)] задача о назначениях является задачей минимизации,
        \item [б)] задача о назначениях является задачей максимизации.
    \end{itemize}
\end{enumerate}

В программе необходимо предусмотреть два режима работы: <<итоговый>>, когда программа печатает только матрицу назначений, и <<отладочный>>, когда на экран поитерационно выводится вспомогательная информация о текущей эквивалентной матрице стоимостей, текущей системе независимых нулей, выделенных столбцах и строках, нулях со штрихами, $L$-цепочке и т.д.